\documentclass{beamer}

\newcommand{\transition}[1]{\frame{{} \begin{center}\fbox{#1}\end{center}}}

\title{MATH 1190 \& 1210 Meeting for Week 4}
\author{Josh}
\date{2 February 2026}


\begin{document}

\frame{\titlepage}


\transition{Content: Application Day 1, The Keeling Curve \& Climate Change}

\section{Logistics}

\transition{Logistics}

\frame[t]{{Logistics}

\begin{itemize}
\item lesson pacing
    \begin{itemize}
    \item {\bf Pacing:} Where is your section in the lessons right now?
    \end{itemize}

\item Week 4 Plan
    \begin{itemize}
    \item TuTh: Lesson 6 on Tuesday + Application Day 1 on Thursday
     \item TuThFr: begin lesson 6 on Tuesday + \& begin Application Day 1 on Thursday + finish Application Day 1 on Friday
    \item MoWe: Lesson 5 on Monday + Application Day 1 on Wednesday
    \item MoWeFr: begin lesson 5 on Monday + finish lesson 5 \& begin Application Day 1 on Wednesday + finish Application Day 1 on Friday
    \end{itemize}
\item Chances to catch up...?
    \begin{itemize}
    \item Lesson 5 may not take 75 minutes
    \item lesson 6 is 2 pages of the power rule, sum rule, difference rule, constant multiple rule, etc. $\implies$ should not take more than 60 minutes
    \item so, there may be some space to catch up a bit in the next week and a half
    \item \textbf{Note:} KC2 (3/10) and KC3 (3/31) will be held on Tuesdays. So, we need to be mindful of pacing.
    \end{itemize}
\end{itemize}

}

\frame{{Logistics}

\begin{itemize}
\item application day logistics
    \begin{itemize}
    \item graded group activity, ideally done in groups of 3
    \item {\bf instructors must make their own assignments in Gradescope for Application Day 1}
        \begin{itemize}
        \item one submission per group
        \item points must total up to 10; how you decide to allot them is up to you
        \item grading for good-effort completion and providing supportive feedback is highly encouraged
        \item if you want to take up your student's paper copies, then choose instructor upload; if you want students to upload their work to Gradescope, then choose student upload
        \end{itemize}
    \item everyone must participate - should see handwriting from 3 distinct people (if in groups of 3)
    \item must be submitted at the end of class - please do not ask students to work on this outside of class 
    \item \textbf{Note:} Next week we will use our time to grade a few sample problems to help with consistency in our grading.
    \end{itemize}
\end{itemize}

}

\frame{{Logistics}

\begin{itemize}
\item P2L Visits Feb 18 and 19
 \item Active learning signup sheet.
 \item Grading for KC1
    \begin{itemize}
        \item F1: Holden, Michael
        \item L1: Vadim, Susan
        \item L2: Aaratrick, Jim R.
        \item D1: Jim B., Josh
        \item LAs have been approved to help grade will assign after today.
        \item Exams will be held in: 
            \begin{itemize}
                \item Warner 209: MATH 1190 – sections 100, 200, 300, MATH 1210 – sections 200, 500
                \item Chemistry 402: MATH 1210 – sections 300, 400, 700, 800, 900, 930
            \end{itemize}
    \end{itemize}
\end{itemize}

}

\section{Pedagogy: Cognitive Demand}

\transition{Pedagogy: Ted Talk: ``Math Class Needs a Makeover"}





% Slide 1
\begin{frame}{Overview of the Talk}
  \begin{itemize}
    \item Many students come from a curriculum that often emphasizes procedures over sense-making.
    \item Students learn how to execute steps, but not always how to decide \emph{which} steps matter.
    \item Meyer challenges us to rethink how math problems are posed, not just how they are solved.
  \end{itemize}
\end{frame}


% Slide 2
\begin{frame}{The Problem with Traditional Math Class}
  \begin{itemize}
    \item Many students expect math problems to be fully scaffolded.
    \item This encourages a “paint-by-numbers” mindset.
    \item Common consequences:
    \begin{itemize}
      \item Low perseverance when problems look unfamiliar
      \item Difficulty transferring skills to new contexts
      \item Anxiety around word problems
      \item Formula-driven thinking instead of reasoning
    \end{itemize}
  \end{itemize}
  
  \begin{block}{Discussion Questions}
    \begin{itemize}
      \item Where do you see this mindset show up most clearly in your classes?
      \item Do strong procedural students struggle when structure is removed?
    \end{itemize}
  \end{block}
\end{frame}


% Slide 3
\begin{frame}{What’s Missing: Cognitive Demand}
  \begin{itemize}
    \item Meyer argues that many math tasks have \emph{low cognitive demand}.
    \item Students are told:
    \begin{itemize}
      \item What the question is
      \item What information matters
      \item What method to use
    \end{itemize}
    \item Real problem solving requires deciding all three.
  \end{itemize}
  
  \begin{block}{Discussion Question}
    Which of these decisions do we most often remove for students in calculus?
  \end{block}
\end{frame}


% Slide 4
\begin{frame}{Problem Formulation vs. Problem Solving}
  \begin{itemize}
    \item Meyer emphasizes that \textbf{formulating the problem} is itself mathematical work.
    \item In authentic contexts:
    \begin{itemize}
      \item The question is ambiguous
      \item Some information is missing or irrelevant
      \item Assumptions must be made explicit
    \end{itemize}
    \item Textbook problems often bypass this phase entirely.
  \end{itemize}
  
  \begin{block}{Discussion Questions}
    \begin{itemize}
      \item Where in calculus do students struggle because they never practiced formulation?
      \item What risks do students face if we never let them sit with ambiguity?
    \end{itemize}
  \end{block}
\end{frame}


% Slide 5
\begin{frame}{Three–Act Tasks and Real Math Problems}
  \begin{itemize}
    \item Meyer’s three–act structure:
    \begin{itemize}
      \item Act 1: Present an engaging situation, no math yet
      \item Act 2: Students decide what information they need
      \item Act 3: Resolve and reflect
    \end{itemize}
    \item The mathematical work emerges naturally from curiosity.
  \end{itemize}
  
  \begin{block}{Discussion Question}
    Where could a three–act structure fit naturally in a calculus topic (limits, optimization, rates of change)?
  \end{block}
\end{frame}


% Slide 6
\begin{frame}{The Tension in Coordinated Courses}
  \begin{itemize}
    \item We face real constraints:
    \begin{itemize}
      \item Common exams
      \item Pacing guides
      \item Large enrollments
      \item Varied instructor experience
    \end{itemize}
    \item Meyer’s approach may feel incompatible with coverage demands.
  \end{itemize}
  
  \begin{block}{Discussion Questions}
    \begin{itemize}
      \item Where do you feel the strongest tension between depth and coverage?
      \item What feels realistically changeable—and what does not?
    \end{itemize}
  \end{block}
\end{frame}


% Slide 7
\begin{frame}{Low-Risk Ways to Increase Cognitive Demand}
  \begin{itemize}
    \item Delay formulas until after intuition is discussed
    \item Remove one piece of information from a standard problem
    \item Ask students to generate the question before solving it
    \item Use a short visual prompt instead of a word problem
  \end{itemize}
  
  \begin{block}{Discussion Question}
    Which of these feels like the easiest first step in your own teaching context?
  \end{block}
\end{frame}

% Slide 8
\begin{frame}{For Graduate Teaching Assistants}
  \begin{itemize}
    \item Students often see TAs and LAs as models for “how to do math.”
    \item How we respond to uncertainty shapes classroom norms.
  \end{itemize}
  
  \begin{block}{Discussion Questions}
    \begin{itemize}
      \item How do you respond when a student asks, “Is this the right way?”
      \item How can you all support productive struggle without feeling unprepared?
    \end{itemize}
  \end{block}
\end{frame}

% Slide 9
\begin{frame}{Final Reflection}
  \textbf{Big-Picture Discussion Questions}
  \begin{enumerate}
    \item What aspects of Meyer’s vision resonate most strongly with your experience?
    \item Where do you remain skeptical—and why?
    \item What is one concrete change you could try in the next two weeks?
    \item How can we support consistency while allowing pedagogical experimentation?
  \end{enumerate}
\end{frame}







\end{document}
\frame{{Pedagogy: Cognitive Demand \& Growth-Based}

How can we view difficulty of a task from a students' perspective?\\[0.1in]

\textbf{\underline{Cognitive Demand}} 

({\footnotesize adapted from {\color{blue}\href{https://journals.sagepub.com/doi/abs/10.3102/00028312033002455}{Stein et al}}, 1996})

\vfill

$\text{low-level demand }\begin{cases} \text{(1) Memorization tasks}\\
\text{(2) Procedural tasks}\end{cases}$

\vfill

$\text{high-level demand }\begin{cases} \text{(3) Connections tasks}\\
\text{(4) Mathematics tasks}\end{cases}$

\vfill

$\text{no demand }\begin{cases} \text{(0) student disengaged}\end{cases}$

\vfill

A student can experience very different CD levels while engaging in the same task depending on how the task is scaffolded.

}

\frame[t]{{Pedagogy: Cognitive Demand \& Growth-Based}

After practicing how to find derivatives using the power, constant multiple, sum, difference, and natural exponential rules, what would the cognitive demand levels of the following prompts be?\

\vfill

{\bf Example 1.} Fill in the blank to complete the rule $\dfrac{d}{dx}[x^n] = \rule[-8pt]{0.5in}{0.4pt}$

\vfill

{\bf Example 2.} $\dfrac{d}{dx}{\frac13x^3-\frac52x^2+5x+11}$

\vfill

{\bf Example 3.} Find all points on the graph of $f(x)=\frac13x^3-\frac52x^2+5x+11$ at which the tangent line to the graph has a slope of $-1$.

\vfill

}

\frame{{Pedagogy: Cognitive Demand \& Growth-Based}

Let's talk about the proficiency \& cognitive demand.

\

First, recall from our last discussion of proficiency, that {\bf slips} (careless mistakes) should {\it not} impact students' scores and that {\bf errors} (processing mistakes) {\it should} impact students' scores when the errors are relevant to the learning target.

\

{\bf Question:} How should cognitive demand of a task impact our expectations for students?

}

\transition{Content}

\section{Content}

\frame{{Content}

{\bf Derivatives of Basic Functions}
\begin{itemize}
\item algebra review + basic differentiation rules + optional extra practice = short lesson
\item pre-class assignment
    \begin{itemize}
    \item introduces basic differentiation rules with examples
    \end{itemize}
\item in-class lesson
    \begin{itemize}
    \item Warm-up Exercise 1: roots as powers, powers as roots
    \item Warm-up Exercise 2: recall basic rules
    \item Exercise A: try basic rules + roots as powers
    \item Exercise B: try basic rules + recall interpretation of derivative
    \item Exercise C: linear approximation (not a learning target, but related to important idea in 1220)
    \item Extra Practice 1: scaffolds optimization (is a learning target later this semester)
    \end{itemize}
\end{itemize}

}

\end{document}

%%%%%%%%%%%%%%%%%%%%%%

\frame{{Content}

{\bf Product \& Quotient Rules}
\begin{itemize}
\item pre-class assignment
    \begin{itemize}
    \item introduces the product \& quotient rules with an example of each
    \end{itemize}
\item in-class lesson
    \begin{itemize}
    \item Warm-up Exercise 1/Example 1/Exercise A: state the product rule, see the product rule, try the product rule
    \item Warm-up Exercise 2/Example 2/Exercise B: state the quotient rule, see the quotient rule, try the quotient rule, try the quotient \& product rules
    \item Extra Practice 1: derivative as slope + product rule + quotient rule
    \item Extra Practice 2: quotient rule + units of derivative + derivative as IROC + graph of derivative = great {\bf D1} review for assessment 1
    \end{itemize}
\end{itemize}

}

\section{Pedagogy: Research in Motivational Psychology}

\transition{Pedagogy: Research in Motivational Psychology}

\frame[t]{{Pedagogy: Research in Motivational Psychology}

Please take a moment to read the introduction of {\it Expectancy-Value-Cost Model of Motivation} up to (but not including) "Overview of Expectancy-Value Models".

\

\[ M = f(E,V,C) \] 

\

Based on your experience, how much do you think the student archetypes - Rory, Jeff, and Jessica - represent the students in a typical calculus classroom?

\

This semester, we will implement a series of ``Application Day" activities to attend to the motivational needs of {\bf Jeff} - the student who did not see the utility of learning calculus.

}
